% Online Appendix to Section 8 (sec:strategic). Numbers: paper/results/strategic.json and
% si_fixed_inputs.json (final rerun of 2026-10-12). Specification: R/categories/policy/
% iea_policy_index.R, cat_policy_index.R, R/indices/build_policy_index.R, R/indices/build_indices.R.
\section{The Policy index and the Strategic Index: further detail}\label{oa:strategic}

\paragraph{Tagging and strength of IEA policies.} Each policy of the IEA Policies and Measures
database \citep{iea2026pams} is assigned to technologies first from its own technology list and,
failing that, from keyword rules applied to its title and text; a policy matched to no technology
is cross-cutting and counted in every technology. Stages are assigned by keyword rules on the text:
Upstream on whole words for mining, mines, exploration, extraction, ores, leases, concessions and
drilling; Midstream on processing, refining, smelting, manufacturing, factories, plants, components,
supply chains and domestic content; Downstream on deployment, installation, consumers, end use,
adoption, efficiency, demand, charging, grid connection, power plants and generation. A policy that
matches no stage is cross-cutting across the three. Its strength is
$\sigma_h=w_{\mathrm{type}}\,w_{\mathrm{status}}\,w_{\mathrm{juris}}\,w_{\mathrm{tech}}\,w_{\mathrm{stage}}$
with type weights 0.85 (financial incentives), 0.70 (public procurement), 0.50 (regulations and
standards; research, development and demonstration), 0.20 (information and labelling; other) and
0.10 (voluntary and coordination measures); status weights 1 (in force), 0.60 (planned), 0.50
(proposed), 0.20 (expired or unknown) and 0 (ended); jurisdiction weights 1 (national), 0.66
(international), 0.50 (subnational, state, provincial or regional), 0.30 (unknown) and 0.25 (local);
a technology discount of 0.70 for a cross-cutting policy and 0.30 for a technology found only in the
text; and a stage discount of 0.75 for a cross-cutting stage. A policy counted in several cells
carries its full strength in each. Upstream words are matched as whole words, because as substrings
``ore'' and ``mine'' match ``more'', ``Korea'' and ``determine''.

\paragraph{The Climate Action Tracker component.} The Tracker rates 65 economies in the final build:
39 in their own right and 26 EU members through the Union's rating. Forty are rated insufficient,
nine critically insufficient, six highly insufficient and ten almost sufficient; among the 46 rated
coalition members the counts are 31, 7, 4 and 4. Because the component varies across a country's
cells only through $g_t$, it tilts the within-country policy ranking. Within the Strategic Index
sample the median within-country Spearman correlation between $\POL^{*}_{its}$ and $g_t$ is 0.82
for the 40 countries rated insufficient and 0.86 for those rated almost sufficient, $-0.92$ for the
nine rated critically insufficient, and 0.12 for the 103 countries scored on the IEA component alone.
The IEA component itself is uncorrelated with $g_t$ within a country (median $-0.07$ for the 65
countries rated by the Tracker). The mean Policy index is 0.64 for the economies scored on both
components and 0.39 for those scored on the IEA component alone; the IEA component distinguishes
a median of 31 values among the 33 cells of the first group and 11 among those of the second. The
index enters the Strategic Index only through within-country ranks, so these level differences do
not move one country against another.

\paragraph{Technology weight.} The lifecycle medians are 820 gCO$_2$e/kWh for pulverised coal, 490
for gas combined cycle, 48 and 41 for utility and rooftop photovoltaics, 27 for concentrated solar,
38 for geothermal, 12 for nuclear and 11 and 12 for onshore and offshore wind
\citep{schlomer2014annex}; oil-fired generation takes the SRREN median of 840 \citep{moomaw2011annex}
in place of the 720 of the input file, because AR5 reports no oil row. Technologies without a row
take proxies: batteries the solar value, the grid the mean of solar and wind, green hydrogen the mean
of gas and solar, and electric vehicles the mean of six generation technologies. The resulting
$g_t$ is 0 (oil), 0.024 (coal), 0.422 (gas), 0.569 (electric vehicles), 0.694 (green hydrogen),
0.967 (solar, batteries and geothermal), 0.983 (grid), 0.999 (nuclear and wind). The sector prior
$0.1\lambda_s+0.1g_t$ spans 0.025 to 0.175; the gap between oil and wind in the technology term
(0.100) is twice the gap between adjacent stages (0.025 to 0.05).

\paragraph{Sample.} Of the 6,918 country--sector rows of the output, 442 belong to countries
without a Policy index and 6,476 to policy-eligible countries. Among the latter, 1,000 fail because
the Economic Opportunity cell is not rankable, 721 because the Energy Security cell is not
rankable, 979 because neither is, and 622 because the Energy Security cell is missing; 3,154 are
scored. Before the Upstream trade baskets were added on 11 October 2026 the index covered 192 countries;
the new Trade category of the Upstream cells is imputed for many small economies and pushes their
cells over the imputed-share threshold, which leaves 174. Sector fixed effects explain 72.6 percent
of the variance, technology fixed effects 67.5 percent, stage fixed effects 4.2 percent and country
fixed effects 4.4 percent. The highest sector means are Solar Upstream (0.634), Batteries Upstream
(0.631), Nuclear Upstream (0.620) and Batteries Midstream (0.618), the lowest the coal and oil
stages (0.410--0.441). Upstream cells have the highest mean need (0.677, against 0.513 Midstream and
0.552 Downstream).

\paragraph{The endowment screen.} An Upstream cell carries a resource base if the technology-level
variable of the Reserves or of the Production category of the Energy Security pillar (for instance
the uranium resources and mine production behind Nuclear Upstream) is positive for the country;
every technology with such a variable is screened. Unquantified USGS holdings count as zero in
these technology roll-ups. Cells without a resource base carry \texttt{resource\_base} FALSE and the
pick exclusion ``no resource base'' in the output; their index values are unchanged. Of the 41
picks the screen changes, the forgone index has median 0.026 (interquartile range 0.016--0.040);
16 picks are unchanged. Every Nuclear Upstream pick under the screen lies in a member with
identified uranium resources (Chile, the Democratic Republic of the Congo, Finland, Gabon, Greece,
Japan, Saudi Arabia and Sweden).

\paragraph{Variants of the Strategic Index.} Each variant recomputes the index over the headline
cells with the endowment screen kept. Within-member Spearman correlation with the headline (median),
identical picks and members whose headline pick is in the variant's top three:
security as strength 0.36, 5 and 12 of 57 (core 0.38, partners 0.24, negative for ten partners);
strength without the fixed weights 0.01, 1 and 9; need without the fixed weights 0.90, 31 and 49;
opportunity and need ranked within the country 0.89, 20 and 49; no policy term 0.88, 30 and 55;
policy from the IEA component alone (the 46 members scored on both) 0.89, 34 and 42; the October
2026 NIPO component added back (50 members) 0.96, 27 and 44. The NIPO variant uses a GTA NIPO
extract of October 2026 (measures in force on 5 October 2026, refreshed on 13 October with the
debiased strength measure), built exactly as the pipeline builds the NIPO Policy Index rows: 77
jurisdictions with EU members individually, 939 rows in the headline technologies, all of them
economies that already have a Policy index; NIPO enters with weight 4 beside IEA PAMS (3) and CAT
(2), renormalised over the components present. The component stored with the February 2026 data
covered only 23 economies. Requiring both policy components (\texttt{OPSI\_POLICY\_REQUIRE\_ALL=1})
leaves 65 eligible countries and 1,666 Strategic Index cells, identical to the headline in value;
46 members keep their pick and the 11 scored on the IEA component alone have none.
